\documentclass[12pt,a4paper]{article}
\usepackage[margin=1in]{geometry}
\usepackage{amsmath,amssymb}
\usepackage{array}
\usepackage{booktabs}
\usepackage{fancyhdr}
\usepackage{enumitem}

\pagestyle{fancy}
\fancyhf{}
\renewcommand{\headrulewidth}{0pt}

\begin{document}

\begin{center}
\begin{tabular}{|c|c|c|c|c|c|c|c|c|c|}
\hline
\multicolumn{10}{|c|}{USN} \\
\hline
& & & & & & & & & \\
\hline
\end{tabular}

\vspace{0.5cm}

{\Large \textbf{NMAM INSTITUTE OF TECHNOLOGY, NITTE}}\\[0.2cm]
{\small \textit{(An Autonomous Institution affiliated to VTU, Belagavi)}}\\[0.3cm]
{\large \textbf{IV Sem B.E. (CSE) Mid Semester Examinations -II, July 2023}}\\[0.2cm]
{\large \textbf{21CS403 -- MICROPROCESSOR AND EMBEDDED SYSTEM}}\\[0.3cm]
\end{center}

\noindent
\begin{tabular}{p{8cm}r}
\textbf{Duration:} 1 Hour & \textbf{Max. Marks:} 20
\end{tabular}

\vspace{0.3cm}
\begin{center}
\textit{Note: Answer any \textbf{One} full question from \textbf{each Unit}.}
\end{center}

\vspace{0.3cm}

\begin{center}
{\large \textbf{Unit -- I}}
\end{center}

\noindent
\begin{tabular}{|p{10.5cm}|c|c|c|c|}
\hline
\textbf{Question} & \textbf{Marks} & \textbf{BT*} & \textbf{CO*} & \textbf{PO*} \\
\hline
\textbf{1. a)} Write the procedure to convert two ASCII values to the equivalent hexadecimal. Write appropriate comments. & 4 & L*2 & 2 & 1 \\
\hline
\textbf{b)} Develop an assembly language program to extract the \textbf{hour} from the system time and display in decimal at row number 8 and column number 20 in the format: HOUR: 10 where 10 is the current hour extracted from the system time. Write appropriate comments. & 6 & L3 & 2 & 2 \\
\hline
\textbf{2. a)} Explain the interrupts required to perform the following tasks by specifying the necessary parameters involved:
\begin{enumerate}[label=\roman*., nosep]
\item Buffered keyboard input
\item Set system date
\item Get cursor position
\item Create sub directory
\end{enumerate} & 4 & L2 & 2 & 1 \\
\hline
\textbf{b)} Develop an assembly language program to read a string from the keyboard and write the same to a file. Display the appropriate error messages. Write appropriate comments. & 6 & L3 & 2 & 2 \\
\hline
\end{tabular}

\vspace{0.5cm}

\begin{center}
{\large \textbf{Unit -- II}}
\end{center}

\noindent
\begin{tabular}{|p{10.5cm}|c|c|c|c|}
\hline
\textbf{Question} & \textbf{Marks} & \textbf{BT*} & \textbf{CO*} & \textbf{PO*} \\
\hline
\textbf{3. a)} Bring out the various differences between RISC and CISC machines. & 5 & L2 & 3 & 1 \\
\hline
\textbf{b)} Develop an embedded C program to display a message using scrolling effect on LCD for a suitable period of time in both directions. Do not use the built-in function for scrolling. & 5 & L3 & 4 & 2 \\
\hline
\textbf{4. a)} Explain the embedded system software in detail. & 5 & L2 & 3 & 1 \\
\hline
\textbf{b)} Develop an embedded C program to perform up-down counter on a 7-segment display using switches. & 5 & L3 & 4 & 2 \\
\hline
\end{tabular}

\vspace{0.5cm}
\begin{center}
\small{Bloom's Taxonomy, L* Level; CO* Course Outcome; PO* Program Outcome}\\[0.3cm]
******************
\end{center}

\end{document}
